% PDF page 272
\source{gilbarg-trudinger:page-0272}


The operator $Q$ is of \textit{divergence form} if there exists a differentiable vector function
$A(x, z, p) = (A^1(x, z, p), \dots, A^n(x, z, p))$ and a scalar function $B(x, z, p)$ such that

\begin{equation}
Q u = \operatorname{div} A(x, u, D u) + B(x, u, D u), \qquad u \in C^2(\Omega);
\tag{10.5}
\end{equation}

that is, in (10.1),

\[
a^{ij}(x, z, p) = \frac{1}{2}\bigl(D_{p_i}A^j(x, z, p) + D_{p_j}A^i(x, z, p)\bigr).
\]

Unlike the case of linear operators, a quasilinear operator with smooth coefficients
is not necessarily expressible in divergence form.

The operator $Q$ is \textit{variational} if it is the Euler-Lagrange operator corresponding
to a multiple integral

\[
\int_{\Omega} F(x, u, D u)\,dx
\]

where $F$ is a differentiable scalar function; that is, $Q$ is of divergence form (10.5) and

\begin{equation}
A^i(x, z, p) = D_{p_i}F(x, z, p), \qquad B(x, z, p) = -D_z F(x, z, p).
\tag{10.6}
\end{equation}

The ellipticity of $Q$ is equivalent to the strict convexity of the function $F$ with
respect to the $p$ variables.

\textit{Examples}

\begin{enumerate}
\item[(i)] $Q u = \Delta u + (\alpha - 2)\,\dfrac{D_i u\,D_j u}{(1 + |D u|^2)}\,D_{ij}u,\qquad \alpha \geq 1.$
\end{enumerate}

Here

\[
\lambda(x, z, p)=
\begin{cases}
1 & \text{if } \alpha \geq 2,\\[6pt]
\dfrac{1+(\alpha-1)|p|^2}{1+|p|^2} & \text{if } \alpha \leq 2,
\end{cases}
\]

\[
\Lambda(x, z, p)=
\begin{cases}
\dfrac{1+(\alpha-1)|p|^2}{1+|p|^2} & \text{if } \alpha \geq 2,\\[6pt]
1 & \text{if } \alpha \leq 2,
\end{cases}
\]

and

\[
\mathcal{E}(x, z, p)=|p|^2(1+(\alpha-1)|p|^2)/(1+|p|^2).
\]

% PDF page 273
\source{gilbarg-trudinger:page-0273}
\pagestyle{empty}

10. Maximum and Comparison Principles

so that $Q$ is elliptic for all $\alpha\ge 1$ and uniformly elliptic only if $\alpha>1$. Writing
\[
Qu=(1+|Du|^2)^{1-\alpha/2}\,\operatorname{div}(1+|Du|^2)^{\alpha/2-1}Du,
\]
we see that $Q$ is equivalent to a divergence form operator, and moreover that $Q$ is
equivalent to the variational operator associated with the integral
\[
\int_{\Omega}(1+|Du|^2)^{\alpha/2}\,dx.
\]
The equation $Qu=0$ coincides with Laplace's equation when $\alpha=2$, and with the
minimal surface equation when $\alpha=1$. For other values of $\alpha$, this equation arises in
the cracking of plates and the modelling of blast furnaces.

(ii) $Qu=\Delta u+\beta D_i u\,D_j u\,D_{ij}u,\quad \beta\ge 0.$

Here
\[
\lambda(x,z,p)=1
\]
\[
\Lambda(x,z,p)=1+\beta|p|^2
\]
\[
\E(x,z,p)=|p|^2(1+\beta|p|^2)
\]
so that $Q$ is elliptic for all $\beta\ge 0$ and uniformly elliptic only when $\beta=0$ (that is
when $Q$ is the Laplacian). For $\beta>0$, $Q$ is equivalent to the variational operator
associated with the integral
\[
\int_{\Omega}\exp\!\left(\frac{\beta}{2}|Du|^2\right)\,dx.
\]
Note that when $\beta\ge 1$, the minimum and maximum eigenvalues of $Q$ are propor-
tional to those in the case $\alpha=1$ in the previous example. However, the existence
results for these operators will turn out to differ substantially due to the different
growth properties of their $\E$ functions.

(iii) \emph{The Equation of Prescribed Mean Curvature}
Let $u\in C^2(\Omega)$ and suppose the graph of $u$ in $\mathbb R^{n+1}$ has mean curvature $H(x)$ at the
point $(x,u(x))$, $x\in\Omega$. (The mean curvature is understood with respect to the
normal direction along which $x_{n+1}$ is increasing). It follows (see Appendix to
Chapter 14) that $u$ satisfies the equation

(10.7) \[ \mathcal{M}u=(1+|Du|^2)\Delta u-D_i uD_j uD_{ij}u=nH(1+|Du|^2)^{3/2}. \]

Here
\[
\lambda(x,z,p)=1,
\]
\[
\Lambda(x,z,p)=1+|p|^2,
\]


% PDF page 274
\source{gilbarg-trudinger:page-0274}

and

\[
\EE(x, z, p) = |p|^2 .
\]

The operator $\MM$ in (10.7) is equivalent to the operator $Q$ in example (i) when $\alpha = 1$.

\medskip

\noindent\textit{(iv) The Equation of Gas Dynamics}

The stationary irrotational flow of an ideal compressible fluid is described by the equation of continuity, $\operatorname{div} (\rho D u)=0$, where $u$ is the velocity potential of the flow and the fluid density $\rho$ satisfies a density--speed relation $\rho=\rho(|Du|)$. In the case of a perfect gas this relation takes the form

\[
\rho = \left( 1 - \frac{\gamma - 1}{2} |Du|^2 \right)^{1/(\gamma - 1)},
\]

where the constant $\gamma$ is the ratio of specific heats of the gas and $\gamma > 1$. The equation satisfied by the velocity potential $u$ is then

\begin{equation}
\tag{10.8}
\Delta u - \frac{D_i u \, D_j u}{1 - \frac{\gamma - 1}{2} |Du|^2} D_{ij} u = 0,
\end{equation}

which has the eigenvalues

\[
\lambda = \frac{1 - \frac{\gamma + 1}{2} |Du|^2}{1 - \frac{\gamma - 1}{2} |Du|^2}, \qquad \Lambda = 1.
\]

Equation (10.8) is elliptic---and the flow is subsonic---when $|Du| < [2/(\gamma + 1)]^{1/2}$, but is hyperbolic when $[2/(\gamma + 1)]^{1/2} < |Du| < [2/(\gamma - 1)]^{1/2}$. We note that (10.8) becomes the minimal surface equation when $\gamma = -1$.

\medskip

\noindent\textit{(v) The Equation of Capillarity}

The equilibrium shape of a liquid surface with constant surface tension in a uniform gravity field is governed by the equation of capillarity,

\begin{equation}
\tag{10.9}
\operatorname{div}\left( \frac{D u}{\sqrt{1 + |Du|^2}} \right) = \kappa u,
\end{equation}

or equivalently,

\[
\MM u = \kappa u (1 + |Du|^2)^{3/2},
\]

where $\MM$ is the operator defined in (10.7), $u$ is the height of the liquid above an undisturbed reference surface, and $\kappa$ is a constant that is positive or negative


% PDF page 275
\source{gilbarg-trudinger:page-0275}

according to whether the gravitational field is acting downward or upward. In the
absence of gravity, the equation is replaced by the constant mean curvature equation, with $H$ equal to a constant in (10.7). The function $\mathcal{E}$ and the eigenvalues $\lambda$, $\Lambda$ are the same as for (10.7). The natural physical boundary condition for the height $u$ in (10.9), when the fluid is constrained by a fixed rigid boundary, is

\[
\frac{\dfrac{\partial u}{\partial \nu}}{\sqrt{1 + |Du|^2}} = \cos \gamma,
\]

where the contact angle $\gamma$ is the angle between the liquid surface and the fixed boundary, measured within the fluid; $\nu$ is the corresponding normal to the fixed boundary.

\bigskip

\noindent\textbf{10.1. The Comparison Principle}

If $L$ is a linear operator satisfying the hypotheses of the weak maximum principle, Corollary 3.2, and if $u, v \in C^0(\overline{\Omega}) \cap C^2(\Omega)$ satisfy the inequalities $Lu \geq Lv$ in $\Omega$, $u \leq v$ on $\partial \Omega$, we have immediately from Corollary 3.2 that $u \leq v$ in $\Omega$. This comparison principle has the following extension to quasilinear operators.

\medskip

\begin{theorem}[Comparison Principle]\label{gilbarg-trudinger.thm.10.1.comparison-principle}\dcref{Theorem 10.1}\group{chapter-10}\source{gilbarg-trudinger:page-0275}
\textit{Let $u, v \in C^0(\overline{\Omega}) \cap C^2(\Omega)$ satisfy $Qu \geq Qv$ in $\Omega$, $u \leq v$ on $\partial \Omega$, where}

\begin{enumerate}
\item[(i)] \textit{the operator $Q$ is locally uniformly elliptic with respect to either $u$ or $v$;}
\item[(ii)] \textit{the coefficients $a^{ij}$ are independent of $z$;}
\item[(iii)] \textit{the coefficient $b$ is non-increasing in $z$ for each $(x,p) \in \Omega \times \mathbb{R}^n$;}
\item[(iv)] \textit{the coefficients $a^{ij}$, $b$ are continuously differentiable with respect to the $p$ variables in $\Omega \times \mathbb{R} \times \mathbb{R}^n$.}
\end{enumerate}

\medskip

\noindent\textit{It then follows that $u \leq v$ in $\Omega$. Furthermore, if $Qu > Qv$ in $\Omega$, $u \leq v$ on $\partial \Omega$ and conditions (i), (ii) and (iii) hold, (but not necessarily (iv)), we have the strict inequality $u < v$ in $\Omega$.}

\medskip

\noindent\textit{Proof.} Let us assume that $Q$ is elliptic with respect to $u$. Then we have

\[
Qu - Qv = a^{ij}(x, \Du)D_{ij}(u-v) + \bigl(a^{ij}(x, \Du) - a^{ij}(x, \Dv)\bigr)D_{ij}v
+ b(x, u, \Du) - b(x, u, \Dv) + b(x, u, \Dv) - b(x, v, \Dv) \geq 0
\]

so that by writing

\[
w = u - v
\]
\[
a^{ij}(x) = a^{ij}(x, \Du)
\]
\[
\bigl[a^{ij}(x, \Du) - a^{ij}(x, \Dv)\bigr]D_{ij}v + b(x, u, \Du) - b(x, u, \Dv) = b^i(x)D_i w
\]

we see that

\[
Lw = a^{ij}(x)D_{ij}w + b^iD_i w \geq 0
\]


% PDF page 276
\source{gilbarg-trudinger:page-0276}

on $\Omega^+ = \{x \in \Omega \mid w(x) > 0\}$ and $w \le 0$ on $\partial \Omega$. Note that the existence of the locally
bounded functions $b^i$ is guaranteed by condition (iv) and the theorem of the mean.
Consequently, using conditions (i) and (iv) we have from Theorem 3.1 that
$w \le 0$ in $\Omega$. If $Qu > Qv$ in $\Omega$, the function $w$ cannot assume a non-negative maximum
in $\Omega$; (see the proof of Theorem 3.1). Hence $w < 0$ in $\Omega$. If $Q$ is elliptic at $v$ the
result follows from the minimum principle for supersolutions. $\square$
\end{theorem}

A uniqueness theorem for the Dirichlet problem for quasilinear elliptic
operators follows immediately from Theorem 10.1.

\medskip

\begin{theorem}[Dirichlet Uniqueness]\label{gilbarg-trudinger.thm.10.2.dirichlet-uniqueness}\dcref{Theorem 10.2}\group{chapter-10}\source{gilbarg-trudinger:page-0276}\uses{gilbarg-trudinger.thm.10.1.comparison-principle}
\textit{Let $u, v \in C^0(\overline{\Omega}) \cap C^2(\Omega)$, satisfy $Qu = Qv$ in $\Omega$, $u = v$ on $\partial \Omega$,
and suppose that conditions (i) to (iv) in Theorem 10.1 hold. Then $u \equiv v$ in $\Omega$.}
\end{theorem}

\medskip

Condition (ii) in the hypotheses of Theorems 10.1 and 10.2 might appear un-
necessarily restrictive. However, we shall show below that the conclusions of
Theorems 10.1 and 10.2 are not generally valid when the principal coefficients
depend on $z$. The comparison principle, Theorem 10.1, will be useful in the establish-
ment of boundary gradient estimates in Chapter 13.
By using the maximum principle for strong solutions (Theorem 9.1) in place of
Corollary 3.2, we see that Theorems 10.1 and 10.2 remain valid when the functions
$u, v \in C^0(\overline{\Omega}) \cap C^1(\Omega) \cap W^{2,n}_{\mathrm{loc}}(\Omega)$.

\medskip

\noindent\textbf{10.2.\ Maximum Principles}

\medskip

Using Theorem 10.1 we can derive the following quasilinear extension of the a priori
bound (Theorem 3.7), which also illustrates the significance of the $\mathcal{E}$ function.

\medskip

\begin{theorem}[A Priori Bound]\label{gilbarg-trudinger.thm.10.3.apriori-bound}\dcref{Theorem 10.3}\group{chapter-10}\source{gilbarg-trudinger:page-0277}
\textit{Let $Q$ be elliptic in $\Omega$ and suppose that there exist non-negative con-
stants $\mu_1$ and $\mu_2$ such that}

\begin{equation}
\tag{10.10}
\frac{b(x,z,p)\ \mathrm{sign}\ z}{\mathcal{E}(x,z,p)} \le \frac{\mu_1 |p| + \mu_2}{|p|^2}
\qquad \forall(x,z,p) \in \Omega \times \mathbb{R} \times \mathbb{R}^n .
\end{equation}

\medskip

\noindent\textit{Then, if $u \in C^0(\overline{\Omega}) \cap C^2(\Omega)$ satisfies $Qu \ge 0(=0)$ in $\Omega$, we have}

\begin{equation}
\tag{10.11}
\sup_{\Omega} u \le \sup_{\partial \Omega} u^+ + C\mu_2
\end{equation}

\noindent where $C = C(\mu_1,\mathrm{diam}\ \Omega)$.

\medskip

\noindent\textit{Proof.}\enspace Let $u \in C^0(\overline{\Omega}) \cap C^2(\Omega)$ and satisfy $Qu \ge 0$ in $\Omega$, and define the operator
$\bar{Q}$ by

\[
\bar{Q}v = a^{ij}(x,u,Du)D_{ij}v + b(x,u,Du).
\]


% PDF page 277
\source{gilbarg-trudinger:page-0277}

\section*{10.2. Maximum Principles}

We choose a comparison function $v$ as in the proof of Theorem 3.7; namely, for
$\mu_2>0$ we set

\[
v(x)=\sup_{\partial\Omega} u^+ + \mu_2\bigl(e^{\alpha d}-e^{\alpha x_1}\bigr),
\]

where $\Omega$ is assumed to lie in the slab, $0<x_1<d$, and $\alpha\ge \mu_1+1$. Then we have in
$\Omega^+=\{x\in\Omega\mid u(x)>0\}$,

\[
\Qbar v=-\mu_2\alpha^2 a^{11}(x,u,Dv)e^{\alpha x_1}+b(x,u,Dv)
\]

\[
\le -\frac{e^{-\alpha x_1}}{\mu_2}\,\E(x,u,Dv)\left(1-\frac{\mu_1}{\alpha}-\frac{e^{-\alpha x_1}}{\alpha^2}\right)
\qquad \text{by (10.10)}
\]

\[
<0\le \Qbar u.
\]

Hence, by Theorem 10.1, we have $u\le v$ in $\Omega$. The result for $\mu_2=0$ follows by letting
$\mu_2$ tend to zero. $\square$

For uniformly elliptic operators, condition (10.10) is equivalent to a condition
of the form

\[
\text{(10.12)}\qquad \frac{b(x,z,p)\,\operatorname{sign} z}{\lambda(x,z,p)}\le \mu_1|p|+\mu_2,
\qquad \forall (x,z,p)\in\Omega\times\R\times\R^n.
\]

An example of a nonuniformly elliptic operator that satisfies (10.10), but not (10.12),
is given by

\[
Qu=\Delta u + D_i u\,D_j u\,D_{ij}u + (1+|Du|^2).
\]

It is clear from the proof of Theorem 10.3 that in the hypotheses we need only have
assumed that (i) $\E>0$ in $\Omega\times\R\times\R^n$; (ii) $Q$ is elliptic with respect to $u$; and (iii)
there exists a fixed vector $p_0\in\R^n$ such that (10.10) hold for all $(x,z,tp_0)$ with
$(x,z,t)\in\Omega\times\R\times\R$. Further maximum principles are treated in Problems 10.1,
10.2.

The condition (10.12) may alternatively be generalized to nonuniformly elliptic
operators by using the Aleksandrov maximum principle (Theorem 9.1). Following
the notation in Chapter 9, we set

\[
\mathscr{D}=\det[a^{ij}(x,z,p)],\qquad \mathscr{D}^*=\mathscr{D}^{1/n}.
\]

\end{theorem}

\begin{theorem}[Aleksandrov-Type Bound]\label{gilbarg-trudinger.thm.10.4.aleksandrov-type-bound}\dcref{Theorem 10.4}\group{chapter-10}\source{gilbarg-trudinger:page-0278}
\textit{Let $Q$ be elliptic in $\Omega$ and suppose there exist non-negative constants
$\mu_1$ and $\mu_2$ such that}

\[
\text{(10.13)}\qquad \frac{b(x,z,p)\,\operatorname{sign} z}{\mathscr{D}^*}\le \mu_1|p|+\mu_2
\qquad \forall (x,z,p)\in\Omega\times\R\times\R^n.
\]


% PDF page 278
\source{gilbarg-trudinger:page-0278}

Then, if $u \in C^0(\overline{\Omega}) \cap C^2(\Omega)$ satisfies $Qu \ge 0$ $(=0)$ in $\Omega$, we have

\begin{equation}
\sup_{\Omega} u(|u|) \le \sup_{\partial\Omega} u^+(|u|) + C\mu_2
\tag{10.14}
\end{equation}

where $C = C(\mu_1, \operatorname{diam}\Omega)$.

\textit{Proof.} In the subdomain $\Omega^+ = \{x \in \Omega \mid u(x) > 0\}$, we have

\[
0 \le Qu = a^{ij}D_{ij}u + b\,\sign u
\]
\[
\le a^{ij}D_{ij}u + [\mu_1(\sign D_i u)D_i u + \mu_2]\mathcal{D}^{*},
\]

and therefore the estimate (10.14) for $\sup_{\Omega} u$ follows by Theorem 9.1. The full
estimate (10.14) is obtained by replacing $u$ with $-u$ in the proof. $\square$

Theorem 10.4 is in fact implicit in the proof of Theorem 9.1. More generally, we
have from Lemma 9.4, the following result.

\begin{theorem}[Integral Bound]\label{gilbarg-trudinger.thm.10.5.integral-bound}\dcref{Theorem 10.5}\group{chapter-10}\source{gilbarg-trudinger:page-0278}
\textit{Let $Q$ be elliptic in the bounded domain $\Omega$ and suppose there exist
non-negative functions $g \in L^n_{\mathrm{loc}}(\mathbb{R}^n)$, $h \in L^n(\Omega)$ such that}

\begin{equation}
\frac{b(x,z,p)\,\sign z}{n\mathcal{D}^{*}} \le \frac{h(x)}{g(p)}
\qquad \forall (x,z,p) \in \Omega \times \mathbb{R} \times \mathbb{R}^n,
\tag{10.15}
\end{equation}

\begin{equation}
\int_{\Omega} h^n\,dx < \int_{\mathbb{R}^n} g^n\,dp = g_{\infty}.
\tag{10.16}
\end{equation}

\textit{Then if $u \in C^0(\overline{\Omega}) \cap C^2(\Omega)$ satisfies $Qu \ge 0$ $(=0)$ in $\Omega$, we have}

\begin{equation}
\sup_{\Omega} u(|u|) \le \sup_{\partial\Omega} u^+(|u|) + C\,\operatorname{diam}\Omega.
\tag{10.17}
\end{equation}

\textit{where $C$ depends on $g$ and $h$.}
\end{theorem}

The quantity $g_{\infty}$, as is the case in Theorem 9.1, may be infinite so that (10.16)
becomes superfluous. If the function $g$ is positive and $G$ is defined by

\[
G^{-1}(t) = \int_{B_t(0)} g^n\,dp
\]

so that $G:(0,g_{\infty}) \to (0,\infty)$, then the constant $C$ in (10.17) is given by

\[
C = G\!\left(\int_{\Omega} h^n\right).
\]
\end{theorem}


% PDF page 279
\source{gilbarg-trudinger:page-0279}




To complete this section we consider the application of Theorem 10.5 to the equation of prescribed mean curvature (10.7). Here

\[
\mathscr{D} = (1 + |p|^2)^{n-1},
\]

so that we may take

\[
g(p) = (1 + |p|^2)^{-(n+2)/2n}.
\]

By a calculation, we obtain

\[
g_\infty = \int_{\R^n} \frac{dp}{(1 + |p|^2)^{n/2 + 1}} = \omega_n,
\]

and hence we have the following estimate:

\begin{corollary}[Prescribed Mean Curvature Bound]\label{gilbarg-trudinger.cor.10.6.prescribed-mean-curvature-bound}\dcref{Corollary 10.6}\group{chapter-10}\source{gilbarg-trudinger:page-0279}\uses{gilbarg-trudinger.thm.10.5.integral-bound}
\textit{Let $u \in C^0(\overline{\Omega}) \cap C^2(\Omega)$ be a solution of the prescribed mean curvature equation (10.7) in the bounded domain $\Omega$. Then, if}

\begin{equation}
H_0 = \int_\Omega |H(x)|^n\,dx < \omega_n,
\tag{10.18}
\end{equation}

\textit{we have}

\begin{equation}
\sup_\Omega |u| \le \sup_{\partial\Omega} |u| + C\,\diam \Omega,
\tag{10.19}
\end{equation}

\textit{where $C = C(n, H_0)$.}
\end{corollary}

Finally we remark that the estimates of this section continue to be valid for subsolutions or solutions $u$ assumed only in $C^0(\overline{\Omega}) \cap W^{2,n}_{\mathrm{loc}}(\Omega)$.

\section*{10.3. A Counterexample}

The following example shows that Theorems 10.1 and 10.2 cannot in general be extended to allow the principal coefficients $a^{ij}$ to depend on $u$. We consider an operator $\mathcal{Q}$ of the form

\begin{equation}
\mathcal{Q}u = \Delta u + g(r,u)\,\frac{x_i x_j}{r^2} D_{ij}u,\qquad r = |x|.
\tag{10.20}
\end{equation}

in the spherical shell $\Omega = \{x \in \R^n \mid 1 < |x| < 2\}$. If $u = u(r)$, the equation $\mathcal{Q}u = 0$ is equivalent to the ordinary differential equation

\[
u'' + u'\left(\frac{n-1}{r(1+g)}\right) = 0.
\]


% PDF page 280
\source{gilbarg-trudinger:page-0280}

Let $v$ and $w$ be polynomials satisfying the conditions

\begin{enumerate}
\item[(i)] $v(1)=w(1), \quad v(2)=w(2);$
\item[(ii)] $v',\, w' >0 \quad \text{in } [1,2];$
\item[(iii)] $v'(1)<w'(1), \quad v'(2)>w'(2);$
\item[(iv)] $v'',\, w'' <0 \quad \text{in } [1,2];$
\item[(v)] $\dfrac{w''(1)}{w'(1)}=\dfrac{v''(1)}{v'(1)}, \quad \dfrac{w''(2)}{w'(2)}=\dfrac{v''(2)}{v'(2)};$
\end{enumerate}

and define for $1\le r\le 2$, $v\le u\le w,$

\[
f(r,u)=\frac{u-v}{w-v}\left(\frac{v''}{v'}-\frac{w''}{w'}\right)-\frac{v''}{v'},
\]

\[
g(r,u)=-1+\frac{n-1}{r f(r,u)}.
\]

Then, writing $v(x)=v(|x|)$, $w(x)=w(|x|)$, we see that $Qv=Qw=0$ in $\Omega$ and $v=w$
on $\partial\Omega$. Also $Q$ is elliptic with respect to both $v$ and $w$. Furthermore, by extending
$f$ in an appropriate way to the strip $[1,2]\times \mathbb{R}$, we can obtain an operator $Q$
that is uniformly elliptic in $\Omega$ and whose coefficients belong to $C^\infty(\overline{\Omega}\times\mathbb{R})$.

\bigskip

\noindent\textbf{10.4. Comparison Principles for Divergence Form Operators}

\medskip

Interesting variants of Theorem 10.1 can be obtained when the operator $Q$ is of
divergence form (10.5). We recall from Chapter 8 that a function $u$, weakly differen-
tiable in $\Omega$, satisfies $Qu\ge 0$ ($=0,\le 0$) in $\Omega$ if the functions $A^i(x,u,Du)$, $B(x,u,Du)$
are locally integrable in $\Omega$ and

\begin{equation}
\tag{10.21}
Q(u,\varphi)=\int_{\Omega}\bigl(A(x,u,Du)\cdot D\varphi-B(x,u,Du)\varphi\bigr)\,dx\le 0\ (=0,\ \ge 0)
\end{equation}

for all non-negative $\varphi\in C_0^1(\Omega)$. The following theorem provides three alternative
criteria for a comparison principle.

\medskip

\begin{theorem}[Comparison for Divergence Form]\label{gilbarg-trudinger.thm.10.7.comparison-divergence-form}\dcref{Theorem 10.7}\group{chapter-10}\source{gilbarg-trudinger:page-0281}
\textit{Let $u,v\in C^1(\overline{\Omega})$ satisfy $Qu\ge 0$ in $\Omega$, $Qv\le 0$ in $\Omega$ and $u\le v$ on $\partial\Omega$,
where the functions $A$, $B$ are continuously differentiable with respect to the $z, p$
variables in $\overline{\Omega}\times\mathbb{R}\times\mathbb{R}^n$, the operator $Q$ is elliptic in $\Omega$, and the function $B$ is non-
increasing in $z$ for fixed $(x,p)\in\Omega\times\mathbb{R}^n$. Then, if either}

\begin{enumerate}
\item[(i)] \textit{the vector function $A$ is independent of $z$; or}
\item[(ii)] \textit{the function $B$ is independent of $p$; or}
\item[(iii)] \textit{the $(n+1)\times(n+1)$ matrix,}
\end{enumerate}

\[
\begin{bmatrix}
D_{p_j}A^i(x,z,p) & -D_{p_j}B(x,z,p)\\
D_z A^i(x,z,p) & -D_z B(x,z,p)
\end{bmatrix}\ge 0 \text{ in } \Omega\times\mathbb{R}\times\mathbb{R}^n;
\]

\textit{it follows that $u\le v$ in $\Omega$.}


% PDF page 281
\source{gilbarg-trudinger:page-0281}

\begin{proof}
Let us define

\[
w=u-v,
\]
\[
u_t=tu+(1-t)v,\qquad 0\le t\le 1,
\]
\[
a^{ij}(x)=\int_0^1 D_{p_j}A^i(x,u_t,Du_t)\,dt,
\]
\[
b^i(x)=\int_0^1 D_zA^i(x,u_t,Du_t)\,dt,
\]
\[
c^i(x)=\int_0^1 D_{p_i}B(x,u_t,Du_t)\,dt,
\]
\[
d(x)=\int_0^1 D_zB(x,u_t,Du_t)\,dt.
\]

Then we have

\begin{align*}
0 &\ge Q(u,\varphi)-Q(v,\varphi) \tag{10.22}\\
  &= \int_\Omega \{((A(x,u,Du)-A(x,v,Dv)))\cdot D\varphi \\
  &\qquad -(B(x,u,Du)-B(x,v,Dv))\varphi\}\,dx \\
  &= \int_\Omega \{(a^{ij}(x)D_jw+b^i(x)w)D_i\varphi-(c^i(x)D_iw+d(x)w)\varphi\}\,dx
\end{align*}

for all non-negative $\varphi\in C^1_0(\Omega)$. Hence $Lw\ge 0$ where $L$ is the linear operator given by

\[
Lw=D_i(a^{ij}D_jw+b^iw)+c^iD_iw+dw.
\]

Since $u,v\in C^1(\overline{\Omega})$, there exist by the hypotheses positive constants $\lambda,\Lambda$ such that

\[
a^{ij}(x)\xi_i\xi_j\ge \lambda |\xi|^2\qquad \forall \xi\in \mathbb R^n,\ x\in \Omega.
\]
\[
|a^{ij}|,|b^i|,|c^i|,|d|\le \Lambda \quad \text{in } \Omega,\qquad d\le 0 \text{ in } \Omega,
\]

and therefore $L$ is strictly elliptic in $\Omega$ with bounded coefficients. The conclusion of Theorem 10.7 can now be obtained directly from the theory of Chapter 8. In particular if condition (i) holds, then $b^i=0$ in $\Omega$ so that by the weak maximum principle, Theorem 8.1, we have $w\le 0$ in $\Omega$. Although the remainder of Theorem 10.7 follows immediately from Problem 8.1, we include the full proof here. If condition (ii)
\end{proof}
\end{theorem}

% PDF page 282
\source{gilbarg-trudinger:page-0282}

holds, then $c^i = 0$ in $\Omega$. Note that this condition on $L$ is equivalent to the preceding
condition on the adjoint operator $L^*$. For $\epsilon > 0$, we define

\[
\varphi = \frac{w^+}{w^+ + \epsilon} \in W^{1,2}_0(\Omega),
\]

and obtain from substitution in (10.22),

\[
\lambda \int_{\Omega} \left|D \log \left(1 + \frac{w^+}{\epsilon}\right)\right|^2\,dx
\le \int_{\Omega} \frac{a^{ij}(x)D_i w^+ D_j w^+}{(w^+ + \epsilon)^2}\,dx
\]
\[
\le \Lambda \int_{\Omega} \frac{w^+}{w^+ + \epsilon}\left|D \log \left(1 + \frac{w^+}{\epsilon}\right)\right|\,dx.
\]

Hence using Young's inequality (7.6), we have

\[
\int_{\Omega} \left|D \log \left(1 + \frac{w^+}{\epsilon}\right)\right|^2\,dx
\le \left(\frac{\Lambda}{\lambda}\right)^2 |\Omega|;
\]

it follows from Poincar\'e's inequality (7.44) that

\[
\int_{\Omega} \left|\log \left(1 + \frac{w^+}{\epsilon}\right)\right|^2\,dx
\le C(n,\lambda,\Lambda,|\Omega|).
\]

Letting $\epsilon \to 0$, we see that $w^+$ must vanish in $\Omega$, that is $w \le 0$ in $\Omega$.
Finally, if condition (iii) holds, we choose $\varphi = w^+$ in $\Omega$ and obtain on substitu-

\[
a^{ij}D_i w^+ D_j w^+ + (b^i - c^i)w^+ D_i w^+ - d(w^+)^2 = 0 \quad \text{in } \Omega,
\]

so that, by Young's inequality (7.6),

\[
|Dw^+|^2 \le n\left(\frac{2\Lambda}{\lambda}\right)^2 |w^+|^2 \quad \text{in } \Omega.
\]

Hence, for any $\epsilon > 0$, we have

\[
\left|D \log \left(1 + \frac{w^+}{\epsilon}\right)\right|
\le \frac{2\sqrt{n}\Lambda}{\lambda}\frac{w^+}{w^+ + \epsilon}
\le \frac{2\sqrt{n}\Lambda}{\lambda}.
\]


% PDF page 283
\source{gilbarg-trudinger:page-0283}

and, since $w^+=0$ on $\partial\Omega$, it follows that
\[
\left|\log\left(1+\frac{w^+}{\varepsilon}\right)\right|\le \frac{2\sqrt{n}\,\Lambda}{\lambda}\,\mathrm{diam}\,\Omega.
\]
Letting $\varepsilon\to 0$ we have, as before, $w^+=0$ in $\Omega$, and therefore $w\le 0$ in $\Omega$. $\square$

Note that when condition (i) holds in Theorem 10.7, we need only assume that $u,\ v\in C^0(\overline{\Omega})\cap C^1(\Omega)$ and that the derivatives of the coefficients belong to $C^0(\Omega\times\mathbb{R}\times\mathbb{R}^n)$. This is easily seen by applying the result of Theorem 10.7 to a subdomain $\Omega'\Subset\Omega$. A similar extension is valid in the other cases provided the coefficients satisfy an appropriate uniform structure condition.

\section*{10.5. Maximum Principles for Divergence Form Operators}

When the operator $Q$ is of divergence form, we can derive maximum principles under different hypotheses from those of Theorems 10.3 and 10.4. We shall assume that the functions $A$ and $B$ in (10.5) satisfy the following structure conditions:

For all $(x,z,p)\in\Omega\times\mathbb{R}\times\mathbb{R}^n$ and some $\alpha\ge 1$,

\begin{equation}
\tag{10.23}
 p\cdot A(x,z,p)\ge |p|^\alpha-|a_1 z|^\alpha-a_2^\alpha,
\end{equation}
\[
B(x,z,p)\,\mathrm{sign}\,z\le
\begin{cases}
 b_0|p|^{\alpha-1}+|b_1 z|^{\alpha-1}+b_2^{\alpha-1} & \text{if }\alpha>1,\\
 b_0 & \text{if }\alpha=1;
\end{cases}
\]
here $a_1,a_2,b_0,b_1,b_2$ are non-negative constants. The first inequality in (10.23) can be viewed as a weak ellipticity condition (see Problem 10.3). The development below is similar to the derivation of global estimates for weak solutions of linear elliptic equations in Chapter 8.

\begin{lemma}[Divergence Form Maximum Bound]\label{gilbarg-trudinger.lem.10.8.divergence-form-maximum-bound}\dcref{Lemma 10.8}\group{chapter-10}\source{gilbarg-trudinger:page-0283}
\textit{Let $u\in C^0(\overline{\Omega})\cap C^1(\Omega)$ satisfy $Qu\ge 0$ in $\Omega$, and suppose that $Q$ satisfies the structure conditions (10.23). Then we have}

\begin{equation}
\tag{10.24}
\sup_{\Omega}u\le C\{\|u^+\|_\alpha+(a_1+b_1)\sup_{\partial\Omega}u^++a_2+b_2\}+\sup_{\partial\Omega}u^+
\end{equation}

\textit{where $C=C(n,\alpha,a_1,b_0,b_1,|\Omega|)$.}
\end{lemma}

\textit{Proof.} Let us assume initially that $u\in C^1(\overline{\Omega})$ and $u\le 0$ on $\partial\Omega$ so that $\sup_{\partial\Omega}u^+=0$. The proof follows that of Theorem 8.15 with this difference: in the

% PDF page 284
\source{gilbarg-trudinger:page-0284}

present case, $u$ is assumed bounded at the outset and hence it is unnecessary to
truncate the power functions used as test functions. By writing

\[
k=a_2+b_2,\qquad \bar z=|z|+k,\qquad \bar b=a_1^{\alpha}+b_0^{\alpha}+b_1^{\alpha-1}+1,
\]

we obtain from inequalities (10.23), with the help of Young’s inequality,

\begin{equation}\tag{10.25}
p\cdot A(x,z,p)\ge |p|^{\alpha}-\bar b|\bar z|^{\alpha}
\end{equation}
\[
\bar z\,B(x,z,p)\,\sign z\le
\begin{cases}
\bigl(\mu |p|^{\alpha}+(\mu^{1-\alpha}+1)\bar b\,\bar z^{\alpha}\bigr) & \text{if $\alpha>1$,}\\
\bar b\,\bar z & \text{if $\alpha=1$,}
\end{cases}
\]

where $\mu>0$. Hence substituting in the integral inequality (10.21) the function

\[
\varphi=w^{\beta}-k^{\beta}
\]

where $w=\bar u^{+}=u^{+}+k$, $\beta\ge 1$, and choosing $\mu=\beta/2$, we obtain

\[
\int_{\Omega} w^{\beta-1}|Dw|^{\alpha}\,dx\le C\bar b \int_{\Omega} w^{\alpha+\beta-1}\,dx,
\]

where $C=C(\beta)$. By the Sobolev inequality (7.26), there exists a number $s>\alpha$ such
that

\[
\|w^{\beta}-k^{\beta}\|_{s}\le Cr\left(\int_{\Omega} w^{\beta-1}|Dw|^{\alpha}\,dx\right)^{1/\alpha}
\]

where $r=(\alpha+\beta-1)/\alpha$ and $C=C(n,s,|\Omega|)$. Consequently we have

\[
\|w\|_{rs}\le (Cr)^{1/r}(\bar b)^{1/\alpha r}\|w\|_{r\alpha}
\]

for all $r\ge 1$, and the estimate (10.24) with $\sup_{\partial\Omega} u^{+}=0$ now follows by the itera-
tion argument of Theorem 8.15. To dispense with the initial assumptions con-
cerning $u$, we replace $u$ by $u-L$ where $L=\sup_{\partial\Omega}u^{+}$, and approximate $\Omega$ by domains
$\Omega'\Subset\Omega$. $\square$

Using Lemma 10.8, we can now derive the following apriori estimates for sub-
solutions and solutions of the equation $Qu=0$.

\begin{theorem}[Divergence Form A Priori Estimate]\label{gilbarg-trudinger.thm.10.9.divergence-form-apriori-estimate}\dcref{Theorem 10.9}\group{chapter-10}\source{gilbarg-trudinger:page-0283}\uses{gilbarg-trudinger.lem.10.8.divergence-form-maximum-bound}
\textit{Let $u\in C^{0}(\overline{\Omega})\cap C^{1}(\Omega)$ satisfy $Qu\ge 0\ ( =0)$ in $\Omega$ and suppose that
$Q$ satisfies the structure conditions (10.23) with $\alpha>1$, $b_1=0$ and either $b_0$ or $a_1=0$.
Then we have the estimate}

\begin{equation}\tag{10.26}
\sup_{\Omega}u(|u|)\le C\bigl(a_2+b_2+a_1\sup_{\partial\Omega}u^{+}(|u|)\bigr)+\sup_{\partial\Omega}u^{+}(|u|)
\end{equation}

\textit{where $C=C(n,\alpha,a_1,b_0,|\Omega|)$.}
\end{theorem}

% PDF page 285
\source{gilbarg-trudinger:page-0285}

Proof. As in the proof of Lemma 10.8 we can assume initially that $u\in C^{1}(\barOmega)$ and
$u\le 0$ on $\partial\Omega$. We also assume that $k>0$. The two cases $b_{0}=0$, $a_{1}=0$ will
be considered separately.

\begin{enumerate}
\item[(i)] Suppose $b_{0}=0$. We then substitute

\[
\varphi=\frac{1}{k^{\alpha-1}}-\frac{1}{w^{\alpha-1}},\qquad w=\bar u^{+},
\]

into the integral inequality (10.21) to obtain

\[
(\alpha-1)\int_{\Omega}\left|\frac{Dw}{w}\right|^{\alpha}\,dx\le \alpha b|\Omega|,
\]

so that

\[
\int_{\Omega}\left|D\log\frac{w}{k}\right|^{\alpha}\,dx\le \frac{\alpha}{\alpha-1}b|\Omega|.
\]

Hence, by Poincar\'e's inequality (7.44),

\[
\int_{\Omega}\left|\log\frac{w}{k}\right|^{\alpha}\,dx\le Cb
\]

where $C=C(n,\alpha,|\Omega|)$. Now writing $M=\sup_{\Omega}w$, we have by the proof of Lemma
10.8

\[
\left(\frac{M}{k}\right)^{\alpha}\le C\int_{\Omega}\left(\frac{w}{k}\right)^{\alpha}\,dx
\]
\[
\le C\left(\frac{M}{k}\right)^{\alpha}\left(\log\frac{M}{k}\right)^{-\alpha}\int_{\Omega}\left(1+\left|\log\frac{w}{k}\right|^{\alpha}\right)\,dx,
\]

so that

\[
\left|\log\frac{M}{k}\right|^{\alpha}\le C\int_{\Omega}\left(1+\left|\log\frac{w}{k}\right|^{\alpha}\right)\,dx\le C.
\]

Hence $M\le Ck$ where $C=C(n,\alpha,a_{1},|\Omega|)$.

\item[(ii)] Suppose $a_{1}=0$. The proof is then similar to that of Theorem 8.16. Writing
again $M=\sup_{\Omega}w$, we substitute

\[
\varphi=\frac{1}{|M-w+k|^{\alpha-1}}-\frac{1}{M^{\alpha-1}}
\]


% PDF page 286
\source{gilbarg-trudinger:page-0286}


into (10.21) to obtain

\[
(\alpha-1)\int_{\Omega}\left|\frac{Dw}{M-w+k}\right|^{\alpha}\,dx
\le b_0\int_{\Omega}\left|\frac{Dw}{M-w+k}\right|^{\alpha-1}\,dx
+\left\{\left(\frac{a_2}{k}\right)^{\alpha}+\left(\frac{b_2}{k}\right)^{\alpha-1}\right\}|\Omega|.
\]

Using Young’s inequality (7.5), we then have

\[
\int_{\Omega}\left|D\log\frac{M}{M-w+k}\right|^{\alpha}\,dx\le Cb|\Omega|
\]

where $C=C(\alpha)$, and hence, by Poincaré’s inequality (7.44),

\[
(10.27)\qquad \int_{\Omega}\left|\log\frac{M}{M-w+k}\right|^{\alpha}\,dx\le Cb,
\]

where $C=C(n,\alpha,|\Omega|)$. To proceed further, we take

\[
\varphi=\frac{\eta}{(M-w+k)^{\alpha-1}}
\]

in (10.21), where $\eta\ge 0$, $\operatorname{supp}\eta\subset \operatorname{supp}u^{+}$ and $\eta\in C_0^1(\Omega)$. We then obtain from the structure conditions (10.23) the inequality

\[
\int_{\Omega}\frac{A\cdot D\eta}{(M-w+k)^{\alpha-1}}\,dx\le \int_{\Omega}\left\{b_0\left|\frac{Dw}{M-w+k}\right|^{\alpha-1}+\alpha\left(\frac{a_2}{k}\right)^{\alpha}+\left(\frac{b_2}{k}\right)^{\alpha-1}\right\}\eta\,dx.
\]

\[
\le \int_{\Omega}\left\{b_0\left|D\log\frac{M}{M-w+k}\right|^{\alpha-1}+\alpha\right\}\eta\,dx.
\]

The function $\bar{w}=\log\,[M/(M-w+k)]$ accordingly satisfies $\Qbar \bar{w}\ge 0$ in $\Omega^{+}=\{x\in \Omega\,|\,u(x)>0\}$, where the operator $\Qbar$ fulfills the structure conditions (10.23) with $a_1=b_1=0$ and $a_2,b_2\le \alpha$. Hence, by Lemma 10.8,

\[
\sup_{\Omega}\bar{w}\le C(\|\bar{w}\|_{x}+1)
\]

\[
\le C(n,\alpha,b_0,|\Omega|)\qquad \text{by (10.27)}
\]

Consequently $M\le Ck$. The case $k=0$ is obtained by letting $k$ tend to zero. By removing the conditions, $u\in C^1(\overline{\Omega})$, $u\le 0$ on $\partial\Omega$, as in the proof of Lemma 10.8, we obtain the estimate (10.26) in each case. \qed

\end{enumerate}

% PDF page 287
\source{gilbarg-trudinger:page-0287}

10.5. Maximum Principles for Divergence Form Operators

As a byproduct of the existence theory for the equation of prescribed mean
curvature in Chapter 16, we shall see that Theorem 10.9 cannot be extended to allow
$\alpha=1$ in its hypotheses. The following estimate, which includes the case $\alpha=1$,
requires that the structure constants $a_1$, $b_0$ and $b_1$ be sufficiently small.

\medskip

\begin{theorem}[Small-Data Divergence Form Bound]\label{gilbarg-trudinger.thm.10.10.small-data-divergence-form-bound}\dcref{Theorem 10.10}\group{chapter-10}\source{gilbarg-trudinger:page-0287}
\textit{Let $u \in C^0(\bar{\Omega}) \cap C^1(\Omega)$ satisfy $Qu \geq 0 (=0)$ in $\Omega$ and suppose
that $Q$ satisfies the structure conditions (10.23). Then there exists a positive constant
$C_0 = C_0(\alpha,n)$ such that if}

\begin{equation}
(a_1^{\alpha} + b_0^{\alpha} + b_1^{\alpha-1})|\Omega|^{\alpha/n} < C_0,
\tag{10.28}
\end{equation}

\textit{we have the estimate}

\begin{equation}
\sup_{\Omega} u^+(|u|) \leq C\{(a_1+b_1)\sup_{\partial\Omega} u^+(|u|) + a_2 + b_2\} + \sup_{\partial\Omega} u^+(|u|).
\tag{10.29}
\end{equation}

\textit{where $C = C(n,\alpha,a_1,b_0,b_1,|\Omega|)$.}
\end{theorem}

\medskip

\noindent\textit{Proof.} By virtue of Lemma 10.8 we need only estimate $\|u^+\|_{\alpha}$. As in the preceding
proofs we assume initially that $u \in C^1(\bar{\Omega})$ and $u \leq 0$ on $\partial\Omega$. Substituting $\varphi = u^+ = u$
into the integral inequality (10.21), we obtain by (10.23)

\begin{align*}
\int_{\Omega} |Du|^{\alpha}\,dx
&\leq \int_{\Omega} \{(a_1^{\alpha} + b_1^{\alpha-1})u^{\alpha} + b_0 u|Du|^{\alpha-1} + a_2^{\alpha} + b_2^{\alpha-1}u\}\,dx \\
&\leq \int_{\Omega} \left\{\left(a_1^{\alpha} + b_1^{\alpha-1} + \frac{b_0^{\alpha}}{\alpha \epsilon^{\alpha-1}}\right)u^{\alpha} + (1-1\alpha)\epsilon|Du|^{\alpha} + a_2^{\alpha} + b_2^{\alpha-1}u\right\}\,dx
\end{align*}

for arbitrary $\epsilon>0$, by inequality (7.6). Taking, in particular, $\epsilon=\alpha^{1/(1-\alpha)}$ for $\alpha\neq 1$,
and using the Poincar\'e inequality (7.44) we obtain for $\alpha\geq 1$

\[
\int_{\Omega} v^{\alpha}\,dx \leq C(n,\alpha)|\Omega|^{\alpha/n}\int_{\Omega}\{(a_1^{\alpha} + b_1^{\alpha-1} + b_0^{\alpha})v^{\alpha} + a_2^{\alpha} + b_2^{\alpha-1}v\}\,dx.
\]

Hence if

\[
C(n,\alpha)|\Omega|^{\alpha/n}(a_1^{\alpha} + b_1^{\alpha-1} + b_0^{\alpha}) < 1,
\]

\textit{we have}

\[
\int_{\Omega} v^{\alpha}\,dx \leq C(a_2^{\alpha}+b_2^{\alpha}),
\]

and the desired estimate (10.29) follows. $\square$


% PDF page 288
\source{gilbarg-trudinger:page-0288}

\begin{quote}
Note that when $\alpha = 1$ in Theorem 10.10, the constants $b_1$ and $b_2$ are not present in inequalities (10.28) and (10.29). By invoking the Poincar\'e inequality (7.44) in the sharp form
\begin{equation}
\tag{10.30}
\int_{\Omega} |v|\,dx \le \frac{1}{n} \left(\frac{|\Omega|}{\omega_n}\right)^{1/n} \int_{\Omega} |Dv|\,dx, \qquad v \in W^{1,1}_0(\Omega),
\end{equation}
we can in this case take
\[
C_0 = C_0(1,n) = n\omega_n^{1/n}.
\]
By writing the equation of prescribed mean curvature (10.7) in its divergence form,
\begin{equation}
\tag{10.31}
\operatorname{div}\,\frac{Du}{\sqrt{1+|Du|^2}} = nH,
\end{equation}
we see that it satisfies the structure conditions (10.23) with $\alpha = 1$, $a_1 = 0$, $a_2 = 1$, $b_2 = n\sup_{\Omega}|H|$. Hence, if the function $H$ satisfies
\begin{equation}
\tag{10.32}
H_0 = \sup_{\Omega}|H| < (\omega_n/|\Omega|)^{1/n},
\end{equation}
we have for any $C^2(\Omega) \cap C^0(\overline{\Omega})$ subsolution (solution) of equation (10.31) the estimate
\begin{equation}
\tag{10.33}
\sup_{\Omega} u(|u|) \le \sup_{\partial\Omega} u(|u|) + C(n,|\Omega|,H_0).
\end{equation}
We conclude this section by noting that the structure conditions (10.23) can be generalized to allow the quantities $a_1$, $a_2$, $b_0$, $b_1$, $b_2$ to be nonnegative measurable functions. In particular if we assume $a_1$, $a_2$, $b_0$, $b_1$, $b_2 \in L^q(\Omega)$ where $q$ is such that $q \ge \alpha$ and $q > n$ and $\bar b_1 = b_1^{1-1/\alpha}$, $\bar b_2 = b_2^{1-1/\alpha}$, then Lemma 10.8 and Theorem 10.9 continue to hold provided in inequalities (10.24) and (10.26), $a_1$, $a_2$, $b_0$, $b_1$, $b_2$ are replaced respectively by $\|a_1\|_q$, $\|a_2\|_q$, $\|b_0\|_q$, $\|\bar b_1\|_{q/(\alpha-1)}$, $\|\bar b_2\|_{q/(\alpha-1)}$ and the constants $C$ depend additionally on $q$. Theorem 10.10 can be similarly extended, the condition (10.28) being replaced by
\begin{equation}
\tag{10.34}
\left\|a_1^{\alpha} + b_0^{\alpha} + b_1^{\alpha-1}\right\|_{\beta} < C_0
\end{equation}
where $\beta = \max(1,n/\alpha)$. For the example of the equation of prescribed mean curvature (10.31) mentioned above, we obtain the more general result (already demonstrated in Corollary 10.6 by different means) that if the function $H$ satisfies
\begin{equation}
\tag{10.35}
\int_{\Omega} |H|^n\,dx < \omega_n,
\end{equation}
\end{quote}

% PDF page 289
\source{gilbarg-trudinger:page-0289}

then the maximum principle (10.33) holds with $H_0 = \|H\|_{n}$. The proofs of these
assertions are basically the same as the case when $a_1, a_2, b_0, b_1, b_2$ are constant; (see
Problem 10.4). Finally we remark that all the results of this section and their proofs
are still applicable if the function $u$ belongs to the Sobolev space $W^{1,1}_{\ast}(\Omega)$ instead of
the space $C^2(\Omega) \cap C^0(\overline{\Omega})$.

\bigskip

\textbf{Notes}

The early results of this chapter, Theorems 10.1, 10.2, and 10.3 are basically
variants of the Hopf maximum principle, Theorem 3.1. The counterexample in
Section 10.3 is due to Meyers [ME 2]. Parts (i) and (ii) of the comparison principle,
Theorem 10.7, were proved in Trudinger [TR 10]. Part (ii) extends an earlier result
of Douglas, Dupont and Serrin [DDS]. Part (iii) was essentially proved in Serrin
[SE 3]. The maximum principle, Theorem 10.9, is a new result in this work although
the technique of its proof has already been demonstrated in [TR 7]. For further
maximum principles of quasilinear equations the reader is referred to the works
[SE 3], [SE 5].

\noindent
We also note here that the form (10.30) of the Poincar\'e inequality is a conse-
quence of the isoperimetric inequality; (see for example [FE]). Theorem 10.5 and
Corollary 10.6 appear in Bakelman [BA 5].

\bigskip

\textbf{Problems}

Use the comparison principle, Theorem 10.1, to establish the following maximum
principles in Problems 10.1, 10.2.

\medskip

\textbf{10.1.} Let $Q$ be elliptic in $\Omega \times \mathbb{R} \times \{0\}$ with coefficients $a^{ij}, b, i,j=1,\dots,n$, differen-
tiable with respect to the $p$ variables in $\Omega \times \mathbb{R} \times \mathbb{R}^n$. Suppose that there exists a
constant $M$ such that

\begin{equation}
\tag{10.36}
zb(x,z,0)\le 0 \qquad \text{for $x\in\Omega$, $|z|\ge M$.}
\end{equation}

Then, if $u\in C^0(\overline{\Omega})\cap C^2(\Omega)$ satisfies $Qu\ge 0\,(=0)$ in $\Omega$, we have

\begin{equation}
\tag{10.37}
\max_{\Omega} u(|u|)\le \max\left\{M,\max_{\partial\Omega} u^+(|u|)\right\}.
\end{equation}

\medskip

\textbf{10.2.} Let $\Omega$ lie in a ball $B_R$ of radius $R$ and suppose that $Q$ is elliptic in $\Omega$ and that

\begin{equation}
\tag{10.38}
(\text{sign } z)\, b(x,z,p)\le \frac{|p|}{R}\,\mathcal{T}(x,z,p), \qquad \mathcal{T}=\trace[a^{ij}].
\end{equation}

for all $x\in\Omega$, $|z|\ge M$, $|p|\ge L$ for constants $M$ and $L$. Then, if $u\in C^0(\overline{\Omega})\cap C^2(\Omega)$
satisfies $Qu\ge 0\,(=0)$ in $\Omega$, we have ([SE 3])

\begin{equation}
\tag{10.39}
\max_{\Omega} u(|u|)\le \max\left\{M,\max_{\partial\Omega} u^+(|u|)\right\}+2LR.
\end{equation}

\noindent
(Hint: Follow the proof of Theorem 10.3 with the half-space, $x_1>0$, replaced by
$B_R$.)

% PDF page 290
\source{gilbarg-trudinger:page-0290}

10.3. Let $Q$ be an operator in the divergence form (10.5). Show that

\begin{equation}\tag{10.40}
p\cdot A(x,z,p)=\int_{0}^{1}s^{-2}\mathcal{G}(x,z,sp)\,ds+p\cdot A(x,z,0).
\end{equation}

Hence show that if $\mathcal{G}\ge c|p|^{\alpha}$, where $c>0$ and $\alpha>1$, then

\begin{equation}\tag{10.41}
p\cdot A(x,z,p)\ge \frac{c}{\alpha-1}|p|^{\alpha}+p\cdot A(x,z,0).
\end{equation}

10.4. Verify the assertions made at the end of Section 10.5.

10.5. Let $\A(p)=[a^{ij}(p)]$ be the coefficient matrix of the minimal surface operator
given by

\[
a^{ij}(p)=(1+|p|^{2})\delta_{ij}-p_{i}p_{j},\qquad p\in \mathbb{R}^{n}.
\]

Verify that $1$ is an eigenvalue of $\A$ with eigenvector $p$ and that $1+|p|^{2}$ is the only
other eigenvalue with eigenspace consisting of vectors orthogonal to $p$.

10.6. Apply Theorem 10.5 to the equations

\[
(1+|Du|^{2})\Delta u-D_{i}uD_{j}uD_{ij}u=nH(x)(1+|Du|)^{s},
\]

where $0\le s<\infty$.

% PDF page 291
\source{gilbarg-trudinger:page-0291}


\chapter{Topological Fixed Point Theorems and Their Application}

In this chapter the solvability of the classical Dirichlet problem for quasilinear equations is reduced to the establishment of certain a priori estimates for solutions. This reduction is achieved through the application of topological fixed point theorems in appropriate function spaces. We shall first formulate a general criterion for solvability and illustrate its application in a situation where the required a priori estimates are readily derived from our previous results. The derivation of these a priori estimates under more general hypotheses will be the major concern of the ensuing chapters.

The fixed point theorems required for the treatment presented here are obtained as infinite dimensional extensions of the Brouwer fixed point theorem, which asserts that a continuous mapping of a closed ball in $\mathbb{R}^n$ into itself has at least one fixed point.

\section{The Schauder Fixed Point Theorem}

The Brouwer fixed point theorem can be extended to infinite dimensional spaces in various ways. We require first the following extension to Banach spaces.

\begin{theorem}[11.1]
Let $\mathcal{G}$ be a compact convex set in a Banach space $\mathcal{B}$ and let $T$ be a continuous mapping of $\mathcal{G}$ into itself. Then $T$ has a fixed point, that is, $Tx = x$ for some $x \in \mathcal{G}$.
\end{theorem}

\begin{proof}
Let $k$ be any positive integer. Since $\mathcal{G}$ is compact, there exists a finite number of points $x_1,\ldots,x_N \in \mathcal{G}$, where $N = N(k)$, such that the balls $B^i = B_{1/k}(x_i)$, $i = 1,\ldots,N$, cover $\mathcal{G}$. Let $\mathcal{G}_k \subset \mathcal{G}$ be the convex hull of $\{x_1,\ldots,x_N\}$, and define the mapping $J_k : \mathcal{G} \to \mathcal{G}_k$ by
\[
J_k x = \frac{\sum \dist(x,\mathcal{G}-B^i)x_i}{\sum \dist(x,\mathcal{G}-B^i)}.
\]
Clearly $J_k$ is continuous and for any $x \in \mathcal{G}$
\end{proof}
